\section{Производные}
	\tikzsetfigurename{FLT_Derivatives_}

Производные для языков предложил Януш Бжозовский в работе~\sidecite{10.1145/321239.321249}.

\begin{definition}[Производная Бжозовского]
    Производная языка $L$ по символу $c$~--- это множество таких цепочек, что приписывание к ним в начало символа $c$ даёт цепочку из языка $L$.
    Записывается производная следующим образом: \[\partial_c L = \{ w' \mid w \in L, w = cw'\}.\]
\end{definition}

Заметим, что если для слова $w$, $|w| = n$ верно, что
\[\varepsilon \in (\partial_{w[n-1]} \circ \dots \circ \partial_{w[1]}  \circ \partial_{w[0]}) (L),\]
то $w \in L$.
Таким образом, существует возможность использовать производные для проверки принадлежность слов заданному языку.
Данная возможность активно используется для регулярных языков~\sidecite{caron_champarnaud_mignot_2014}, языков, распознаваемых автоматами, управляемыми входом (Input Driven Pushdown Automata, IDPDA)~\sidecite{10.1145/3591472}, и контекстно-свободных языков~\sidecite{10.1145/2034773.2034801, 10.1145/3022671.2984026}.
Кроме этого, с помощью производных можно построить элегантный алгоритм построения конечного автомата по регулярному выражению~\sidecite{10.1017/S0956796808007090} и для выполнения регулярных запросов~\sidecite{10.1145/2949689.2949711}.
